\section{Atoms, Stoichiometry, and Nomenclature}

\subsection{Ctrl+F keywords: isotope, neutron, molar mass, empirical formula, limiting reactant, yield, navngivning, stofmaengde}
Use for exam tasks asking for isotope particles, formula mass, mass percent, molecular formula, reaction amounts, limiting reagent, theoretical yield, and ionic compound names.

\subsection{Core formulas and definitions}

\subsubsection{Isotope notation, mass number, atomic number, neutron count}
\[
{}^{A}_{Z}\ce{X}, \qquad A=Z+N, \qquad N=A-Z
\]
\(Z\) is protons, \(N\) is neutrons, and a neutral atom has \(Z\) electrons. Isotopes have the same \(Z\) but different \(N\).

\subsubsection{Mole, Avogadro constant, molar mass}
\[
N=nN_{\mathrm A}, \qquad N_{\mathrm A}=6.022\times10^{23}\,\mathrm{mol^{-1}}, \qquad n=\frac{m}{M}
\]
Convert all masses to moles before using reaction coefficients. \(M\) is found by summing atomic molar masses in the formula.

\subsubsection{Mass percent and composition}
\[
w_i=\frac{m_i}{m_{\mathrm{total}}}\times100\%, \qquad
w_{\mathrm{element}}=\frac{n_iM_i}{M_{\mathrm{compound}}}\times100\%
\]
Use subscripts in the chemical formula as mole counts of atoms per mole compound.

\subsubsection{Empirical formula and molecular formula}
\[
n_i=\frac{m_i}{M_i}, \qquad
\text{empirical ratio}=\frac{n_i}{\min(n_i)}, \qquad
M_{\mathrm{molecular}}=kM_{\mathrm{empirical}}
\]
Round only after checking near-integer ratios such as \(1.5\to3/2\). The molecular formula is \(k\) times the empirical formula.

\subsubsection{Balanced equation stoichiometry}
\[
\nu_A\ce{A}+\nu_B\ce{B}\rightarrow\nu_C\ce{C}, \qquad
n_C=n_A\frac{\nu_C}{\nu_A}
\]
Stoichiometric coefficients are mole ratios, not mass ratios. Balance atoms and charge before calculating.

\subsubsection{Limiting reactant and theoretical yield}
\[
\xi_i=\frac{n_i}{\nu_i}, \qquad
\xi_{\max}=\min(\xi_i), \qquad
n_{\mathrm{product}}=\xi_{\max}\nu_{\mathrm{product}}
\]
The smallest available reaction extent is the limiting reactant. Convert product moles to mass with \(m=nM\).

\subsubsection{Percent yield}
\[
\%\mathrm{yield}=\frac{\mathrm{actual\ yield}}{\mathrm{theoretical\ yield}}\times100\%
\]
Use the same units in numerator and denominator. Percent yield above \(100\%\) usually signals a wrong limiting reactant or unit conversion.

\subsection{Exam recipe: stoichiometry, limiting reactant, gas amount, udbytte}
\textbf{Use when:} the problem gives reactant masses, solution concentrations, or gas data and asks for product mass, amount, excess, or yield.

\textbf{Steps:}
\begin{enumerate}
    \item Write and balance the reaction.
    \item Convert every given reactant to moles.
    \item Divide each reactant amount by its coefficient \(\nu_i\).
    \item The smallest \(n_i/\nu_i\) is limiting.
    \item Multiply \(\xi_{\max}\) by product coefficient and convert to requested units.
\end{enumerate}

\textbf{Fast checks:} atoms and charge balance; answer cannot exceed theoretical yield; units must cancel.

\textbf{Common traps:} using grams as mole ratios; forgetting coefficients for \(\ce{H2}\), \(\ce{O2}\), or polyatomic ions; treating excess reactant as limiting.

\subsection{Exam recipe: empirical formula from mass percent, sum formula}
\textbf{Use when:} percent composition or masses of elements are given.

\textbf{Steps:}
\begin{enumerate}
    \item Assume \(100\,\mathrm g\) sample if percentages are given.
    \item Convert each element mass to moles.
    \item Divide all mole amounts by the smallest.
    \item Multiply all ratios by \(2,3,\ldots\) only if needed to get integers.
    \item Use molar mass to scale empirical to molecular formula if requested.
\end{enumerate}

\textbf{Fast checks:} empirical formula mass must divide the molecular molar mass approximately by an integer.

\subsection{Ionic formulas and names}

\subsubsection{Charge neutrality for ionic compounds}
\[
\sum_i q_in_i=0
\]
Subscripts are chosen so total positive and negative charge cancel. Use parentheses when a polyatomic ion appears more than once, for example \(\ce{(NH4)2HPO4}\).

\subsubsection{Common ions, Danish/English search names}
\[
\begin{array}{c c@{\qquad}c c}
\toprule
\text{ion} & \text{name} & \text{ion} & \text{name}\\
\midrule
\ce{NH4+} & \text{ammonium} & \ce{OH-} & \text{hydroxide}\\
\ce{NO3-} & \text{nitrate} & \ce{NO2-} & \text{nitrite}\\
\ce{SO4^2-} & \text{sulfate/sulfat} & \ce{SO3^2-} & \text{sulfite/sulfit}\\
\ce{HSO4-} & \text{hydrogen sulfate} & \ce{CO3^2-} & \text{carbonate/carbonat}\\
\ce{PO4^3-} & \text{phosphate/fosfat} & \ce{HPO4^2-} & \text{hydrogen phosphate}\\
\ce{H2PO4-} & \text{dihydrogen phosphate} & \ce{CN-} & \text{cyanide}\\
\bottomrule
\end{array}
\]
These names recur in formula writing, buffers, precipitation, and redox questions.

\subsubsection{Binary molecular prefixes}
\[
\begin{array}{c c@{\quad}c c@{\quad}c c}
\toprule
1&\text{mono} & 2&\text{di} & 3&\text{tri}\\
4&\text{tetra} & 5&\text{penta} & 6&\text{hexa}\\
7&\text{hepta} & 8&\text{octa} & 9&\text{nona}\\
10&\text{deca} &&&\\
\bottomrule
\end{array}
\]
The second element receives the ending ``-ide''. Omit ``mono'' on the first element in common naming.
